% =====================================================================
%  main.tex  --  "director de orquesta"
% =====================================================================
%  Este fichero SOLO ensambla el proyecto. No contiene diseño (está en
%  styles/mistilo.sty) ni texto largo (está en capitulos/). Su trabajo:
%    1. Elegir la clase de documento, el idioma y la bibliografía.
%    2. Rellenar los datos del documento (título, autor...).
%    3. Decidir el ORDEN de las partes: portada, índice, secciones,
%       bibliografía.
%
%  CÓMO SE COMPILA (la bibliografía necesita el programa biber):
%     pdflatex main -> biber main -> pdflatex main -> pdflatex main
%  ¿Por qué tantas pasadas? La 1.ª recoge las citas, biber prepara la
%  bibliografía y las dos últimas colocan números y referencias.
%  Más cómodo, con un solo comando:   latexmk -pdf main.tex
%  (En Overleaf basta con pulsar "Recompile".)
%
%  CÓMO LEER UN COMANDO LaTeX
%    \comando[opciones]{argumento}
%    - [opciones] son opcionales; {argumento} es el dato principal.
%    - Todo lo que sigue a un  %  hasta final de línea es comentario.
% =====================================================================

% ---- Clase de documento ----
% "article" sirve para informes y entregables cortos: se divide en
% \section. Para un TFG largo cambia a "report" y usa \chapter.
% 11pt = tamaño base de letra.
\documentclass[11pt]{article}

% ---- Idioma y codificación ----
% inputenc: permite escribir tildes y ñ directamente.
% babel[spanish]: traduce "Contents" a "Índice", pone la separación
% silábica en español, etc.
\usepackage[utf8]{inputenc}
\usepackage[spanish]{babel}

% ---- Bibliografía ----
% biblatex gestiona citas y bibliografía.
%   backend=biber -> el programa que procesa el .bib (necesario).
%   style=ieee    -> formato de las citas: [1], [2]...
% QUÉ PUEDES CAMBIAR: style=apa, style=numeric, style=authoryear...
% Prueba a cambiarlo y recompila: el texto de los capítulos NO se toca,
% solo cambia cómo se ven las citas y la lista final.
\usepackage[backend=biber,style=ieee]{biblatex}
% Indica dónde está la base de datos de referencias.
\addbibresource{bib/refs.bib}

% ---- Nuestro paquete de estilo ----
% Carga todo el diseño. La ruta es relativa a este fichero, así el
% proyecto es portable: se copia la carpeta entera y compila.
\usepackage{styles/mistilo}

% ---- Datos del documento ----
% Para cada entregable nuevo solo se cambian estas líneas.
%   \titulodoc[título corto para la cabecera]{título completo}
\titulodoc[Programación paralela y datos enlazados]{Programación paralela y datos enlazados: un informe de ejemplo}
\autordoc{Nombre Apellido}                         % <- pon tu nombre
\titulacion{Ingeniería informática}                % <- tu grado
\mencion{Tecnologías de la información y la comunicación}   % <- tu mención
% \fechadoc{Octubre de 2026}   % descomenta para fijar la fecha; si no,
%                              % se usa la fecha de compilación

% ---- Compilar solo una sección mientras trabajas ----
% Si el documento es largo, compilar todo tarda. Descomenta esta línea
% para compilar únicamente la sección 2. Con \include (y no \input) las
% demás conservan su numeración. Vuelve a comentarla al terminar.
% \includeonly{capitulos/02_datos}

% =====================================================================
%  A partir de aquí empieza el documento visible
% =====================================================================
\begin{document}

\makeportada                      % portada definida en mistilo.sty

\renewcommand{\contentsname}{Índice}   % título del índice
\tableofcontents                  % genera el índice automáticamente
\newpage

% Cada \include inserta un fichero de capitulos/ SIN escribir ".tex".
% \include empieza en página nueva: ideal para partes independientes.
% (\input insertaría el texto tal cual, útil para fragmentos pequeños.)
% PARA AÑADIR UNA SECCIÓN NUEVA:
%   1) crea capitulos/03_conclusiones.tex con un \section{...}
%   2) añade aquí la línea:  \include{capitulos/03_conclusiones}
% =====================================================================
%  capitulos/01_intro.tex  --  SOLO CONTENIDO
% =====================================================================
%  Este fichero contiene únicamente texto. Fíjate en lo que NO hay:
%    - \documentclass, \usepackage ni \begin{document}: eso es de main.tex
%      y del .sty. main.tex lo incorpora con % =====================================================================
%  capitulos/01_intro.tex  --  SOLO CONTENIDO
% =====================================================================
%  Este fichero contiene únicamente texto. Fíjate en lo que NO hay:
%    - \documentclass, \usepackage ni \begin{document}: eso es de main.tex
%      y del .sty. main.tex lo incorpora con % =====================================================================
%  capitulos/01_intro.tex  --  SOLO CONTENIDO
% =====================================================================
%  Este fichero contiene únicamente texto. Fíjate en lo que NO hay:
%    - \documentclass, \usepackage ni \begin{document}: eso es de main.tex
%      y del .sty. main.tex lo incorpora con \include{capitulos/01_intro}.
%    - colores ni márgenes: eso es diseño y vive en styles/mistilo.sty.
%
%  PISTAS PARA EDITAR
%    - Una línea en blanco separa párrafos.
%    - ~ es un espacio que NO permite salto de línea ("sección~2").
%    - \emph{...} = cursiva de énfasis;  \textbf{...} = negrita.
% =====================================================================

% \section crea una sección numerada y la añade al índice.
\section{Introducción}

La programación paralela es hoy una herramienta cotidiana, pero escribir
código paralelo eficiente y portable sigue siendo difícil. Una línea de
trabajo de la Universidad de Valladolid propone describir los programas
con especificaciones de alto nivel y dejar que un compilador modular
genere el código final~\cite{gonzalez2011trasgo}.

% \cite{clave} inserta una cita. La clave "gonzalez2011trasgo" apunta a
% una entrada de bib/refs.bib. Si la clave no existe o está mal escrita
% saldrá [?] y un aviso al compilar. El formato visible ([1],
% (González, 2011)...) lo decide el estilo elegido en main.tex, no este
% fichero.
%
% Para citar VARIAS referencias a la vez:
%     \cite{gonzalez2011trasgo,fernandez2013hdt}

% Entorno "aviso" definido en mistilo.sty. Sin argumento, el título
% por defecto es "Aviso".
\begin{aviso}
  Este documento es solo un ejemplo para el taller: las referencias son
  reales, el resto es ficticio.
\end{aviso}

% \subsection crea una subsección (1.1, 1.2...).
\subsection{Estructura del informe}
% \ref{etiqueta} pone el número de algo marcado con \label{etiqueta}.
% Aquí "sec:datos" está definida en capitulos/02_datos.tex. Si se
% reordenan las secciones, el número se actualiza solo (tras 2 pasadas).
La sección~\ref{sec:datos} describe un formato binario para publicar
grandes conjuntos de datos enlazados~\cite{fernandez2013hdt} y
presenta una tabla resumen.
.
%    - colores ni márgenes: eso es diseño y vive en styles/mistilo.sty.
%
%  PISTAS PARA EDITAR
%    - Una línea en blanco separa párrafos.
%    - ~ es un espacio que NO permite salto de línea ("sección~2").
%    - \emph{...} = cursiva de énfasis;  \textbf{...} = negrita.
% =====================================================================

% \section crea una sección numerada y la añade al índice.
\section{Introducción}

La programación paralela es hoy una herramienta cotidiana, pero escribir
código paralelo eficiente y portable sigue siendo difícil. Una línea de
trabajo de la Universidad de Valladolid propone describir los programas
con especificaciones de alto nivel y dejar que un compilador modular
genere el código final~\cite{gonzalez2011trasgo}.

% \cite{clave} inserta una cita. La clave "gonzalez2011trasgo" apunta a
% una entrada de bib/refs.bib. Si la clave no existe o está mal escrita
% saldrá [?] y un aviso al compilar. El formato visible ([1],
% (González, 2011)...) lo decide el estilo elegido en main.tex, no este
% fichero.
%
% Para citar VARIAS referencias a la vez:
%     \cite{gonzalez2011trasgo,fernandez2013hdt}

% Entorno "aviso" definido en mistilo.sty. Sin argumento, el título
% por defecto es "Aviso".
\begin{aviso}
  Este documento es solo un ejemplo para el taller: las referencias son
  reales, el resto es ficticio.
\end{aviso}

% \subsection crea una subsección (1.1, 1.2...).
\subsection{Estructura del informe}
% \ref{etiqueta} pone el número de algo marcado con \label{etiqueta}.
% Aquí "sec:datos" está definida en capitulos/02_datos.tex. Si se
% reordenan las secciones, el número se actualiza solo (tras 2 pasadas).
La sección~\ref{sec:datos} describe un formato binario para publicar
grandes conjuntos de datos enlazados~\cite{fernandez2013hdt} y
presenta una tabla resumen.
.
%    - colores ni márgenes: eso es diseño y vive en styles/mistilo.sty.
%
%  PISTAS PARA EDITAR
%    - Una línea en blanco separa párrafos.
%    - ~ es un espacio que NO permite salto de línea ("sección~2").
%    - \emph{...} = cursiva de énfasis;  \textbf{...} = negrita.
% =====================================================================

% \section crea una sección numerada y la añade al índice.
\section{Introducción}

La programación paralela es hoy una herramienta cotidiana, pero escribir
código paralelo eficiente y portable sigue siendo difícil. Una línea de
trabajo de la Universidad de Valladolid propone describir los programas
con especificaciones de alto nivel y dejar que un compilador modular
genere el código final~\cite{gonzalez2011trasgo}.

% \cite{clave} inserta una cita. La clave "gonzalez2011trasgo" apunta a
% una entrada de bib/refs.bib. Si la clave no existe o está mal escrita
% saldrá [?] y un aviso al compilar. El formato visible ([1],
% (González, 2011)...) lo decide el estilo elegido en main.tex, no este
% fichero.
%
% Para citar VARIAS referencias a la vez:
%     \cite{gonzalez2011trasgo,fernandez2013hdt}

% Entorno "aviso" definido en mistilo.sty. Sin argumento, el título
% por defecto es "Aviso".
\begin{aviso}
  Este documento es solo un ejemplo para el taller: las referencias son
  reales, el resto es ficticio.
\end{aviso}

% \subsection crea una subsección (1.1, 1.2...).
\subsection{Estructura del informe}
% \ref{etiqueta} pone el número de algo marcado con \label{etiqueta}.
% Aquí "sec:datos" está definida en capitulos/02_datos.tex. Si se
% reordenan las secciones, el número se actualiza solo (tras 2 pasadas).
La sección~\ref{sec:datos} describe un formato binario para publicar
grandes conjuntos de datos enlazados~\cite{fernandez2013hdt} y
presenta una tabla resumen.

% =====================================================================
%  capitulos/02_datos.tex  --  SOLO CONTENIDO
% =====================================================================
%  Sección 2 del informe. Aquí practicamos: etiquetas y referencias
%  cruzadas, figuras, tablas y cajas de aviso.
% =====================================================================
\section{Datos enlazados a gran escala}
% \label da un "nombre" a esta sección para poder citarla desde otro
% sitio con \ref{sec:datos}. Convención: prefijo + tema
% (sec: sección, fig: figura, tab: tabla).
\label{sec:datos}

El formato \emph{HDT} divide un conjunto de datos RDF en tres
componentes: cabecera, diccionario y tripletas~\cite{fernandez2013hdt}.
Esa separación es, en el fondo, la misma idea que aplicamos hoy a
LaTeX: \textbf{dividir en módulos con una responsabilidad clara}.

% ---- FIGURA ----
% Un entorno "figure" crea una figura numerada con pie y etiqueta.
%   [H]         = colócala exactamente aquí (lo permite el paquete float)
%   \centering  = centrar
%   \includegraphics[width=...]{fichero}: inserta la imagen.
%        width=0.85\linewidth = 85 % del ancho del texto
%        El fichero se busca en Recursos/ gracias a \graphicspath (.sty).
%        Se usa el PDF vectorial (nítido a cualquier zoom), no el PNG.
%   \caption    = texto del pie ("Figura 1: ...")
%   \label      = nombre para citarla con \ref{fig:hdt}
% PARA CAMBIAR LA FIGURA: guarda otra en Recursos/ y cambia el nombre.
\begin{figure}[H]
  \centering
  \includegraphics[width=0.85\linewidth]{hdt_esquema.pdf}
  \caption{Esquema de HDT (figura vectorial).}
  \label{fig:hdt}
\end{figure}

% ---- TABLA ----
% tabular define las columnas con {lll}: tres columnas alineadas a la
% izquierda (c = centrada, r = derecha).
%   &   separa columnas     \\   termina una fila
%   \toprule \midrule \bottomrule: líneas horizontales de booktabs
%   @{}...@{} en la definición quita el margen lateral sobrante.
% PARA AÑADIR UNA FILA: copia una fila existente y cambia su texto
% (no olvides los & ni el \\ final).
\begin{table}[H]
  \centering
  \begin{tabular}{@{}lll@{}}
    \toprule
    \textbf{Componente} & \textbf{Contenido} & \textbf{Función} \\
    \midrule
    Header     & Metadatos        & Describe el conjunto de datos \\
    Dictionary & Cadenas de texto & Evita repetirlas \\
    Triples    & Estructura       & Representa el grafo \\
    \bottomrule
  \end{tabular}
  \caption{Los tres componentes de HDT.}
  \label{tab:hdt}
\end{table}

% \ref{fig:hdt} y \ref{tab:hdt} se sustituyen por los números reales.
La figura~\ref{fig:hdt} y la tabla~\ref{tab:hdt} resumen la idea.

% El mismo entorno "aviso", ahora con título propio (argumento opcional
% entre corchetes).
\begin{aviso}[Regla práctica]
  Cuando copies y pegues un bloque por tercera vez, es hora de
  convertirlo en un entorno.
\end{aviso}


% Imprime la lista de referencias citadas con \cite en el texto.
% Solo aparecen las que realmente has citado.
\printbibliography[heading=bibintoc,title={Bibliografía}]

\end{document}
